% tab:arch-componentes-funcionales
{\small
\setlength{\LTleft}{0pt}
\setlength{\LTright}{0pt}
\setlength{\tabcolsep}{2pt}
\renewcommand{\arraystretch}{1.18}
\begin{longtable}{>{\raggedright\arraybackslash}p{0.30\textwidth} >{\raggedright\arraybackslash}p{0.62\textwidth}}
\caption{Componentes funcionales del prototipo}\label{tab:arch-componentes-funcionales}\\
\hline
\textbf{Componente} & \textbf{Función principal} \\
\hline
\endfirsthead
\caption[]{Componentes funcionales del prototipo (continuaci?n)}\\
\hline
\textbf{Componente} & \textbf{Función principal} \\
\hline
\endhead
\hline
\endfoot
Entrada y preprocesamiento & Cargar el estudio, preservar su información espacial y prepararlo para la inferencia \\
Segmentación & Delimitar las estructuras anatómicas contempladas en cada plano mediante modelos de aprendizaje profundo \\
Mediciones & Calcular valores geométricos derivados de máscaras visibles y trazables \\
Visualización & Mostrar la imagen original con las máscaras superpuestas y sus resultados asociados \\
Salida editable & Organizar los resultados en una plantilla revisable y un preinforme modificable y no diagnóstico \\
Validación & Comparar máscaras contra anotaciones de referencia y complementar con revisión profesional \\
\hline
\end{longtable}
}

% tab:arch-decisiones-requerimientos
{\footnotesize
\setlength{\LTleft}{0pt}
\setlength{\LTright}{0pt}
\setlength{\tabcolsep}{2pt}
\renewcommand{\arraystretch}{1.18}
\phantomsection\label{tab:arch-decisiones-requerimientos}
\begin{longtable}{>{\raggedright\arraybackslash}p{0.24\textwidth} >{\raggedright\arraybackslash}p{0.32\textwidth} >{\raggedright\arraybackslash}p{0.36\textwidth}}
\hline
\textbf{Requerimiento o necesidad} & \textbf{Implicancia técnica} & \textbf{Respuesta arquitectónica} \\
\hline
\endfirsthead
\hline
\textbf{Requerimiento o necesidad} & \textbf{Implicancia técnica} & \textbf{Respuesta arquitectónica} \\
\hline
\endhead
\hline
\endfoot
Carga de estudios de-identificados & El caso debe conservar identidad, plano y metadata sin almacenar datos personales innecesarios & Backend para registrar el estudio y Frontend para seleccionar los volúmenes. \\
Segmentación sagital y axial & Coexisten modelos con datasets, clases y artifacts diferentes & Registro por modelKey, ejecución por plano y agrupación bajo multiplanarRunId. \\
Revisión profesional & La salida automática no puede considerarse definitiva & Estados pendiente, aceptado, observado y descartado, con notas y auditoría. \\
Trazabilidad & Cada corrida debe poder explicarse y reproducirse & caseId, runId, multiplanarRunId, traceId, modelVersion, artifactHash y modos de inferencia. \\
Reapertura sin AI Module & Los resultados deben permanecer disponibles aunque el motor de IA esté apagado & Persistencia durable de contratos, mediciones y assets servidos por el Backend. \\
Navegación volumétrica & Un archivo MHA contiene múltiples cortes y el usuario debe poder recorrerlos & Catálogo de slices, previews persistidos, índice seleccionado y visor navegable. \\
No diagnóstico autónomo & El sistema debe asistir sin reemplazar el juicio profesional & humanReviewRequired, notClinicalDiagnosis y ausencia de conclusiones clínicas automáticas. \\
Actualización de modelos & Los checkpoints pueden evolucionar y exceder los límites de GitHub & Código en repositorios; artifacts y manifests en storage externo verificable. \\
Despliegue controlado & La infraestructura debe ser reproducible y limitar costos & Contenedores, variables de entorno, servicios separados y arquitectura objetivo en Google Cloud. \\
\hline
\end{longtable}
}

% tab:arch-logica-capas
{\footnotesize
\setlength{\LTleft}{0pt}
\setlength{\LTright}{0pt}
\setlength{\tabcolsep}{2pt}
\renewcommand{\arraystretch}{1.18}
\phantomsection\label{tab:arch-logica-capas}
\begin{longtable}{>{\raggedright\arraybackslash}p{0.18\textwidth} >{\raggedright\arraybackslash}p{0.22\textwidth} >{\raggedright\arraybackslash}p{0.52\textwidth}}
\hline
\textbf{Capa} & \textbf{Tecnología} & \textbf{Responsabilidad principal} \\
\hline
\endfirsthead
\hline
\textbf{Capa} & \textbf{Tecnología} & \textbf{Responsabilidad principal} \\
\hline
\endhead
\hline
\endfoot
Frontend & React, Vite y TypeScript & Autenticación visual, worklist, alta y consulta de estudios, workspace sagital–axial, visor de cortes, revisión, historial y diagnósticos técnicos. \\
Backend & Spring Boot y Java & Seguridad, reglas de negocio, persistencia, auditoría, proxy hacia IA, contratos canónicos, serving de assets y reapertura. \\
AI Module & Python y FastAPI & Registry de modelos, manifests, hashing, preprocesamiento, inferencia, reportes, catálogo volumétrico y generación de previews. \\
Runtime de inferencia & PyTorch & Reconstrucción de arquitecturas, carga de state\_dict, model.eval(), torch.inference\_mode() y predicción. \\
Persistencia & PostgreSQL & Usuarios, roles, estudios, corridas, revisiones, mediciones, artifacts y eventos de auditoría. \\
Storage de modelos y archivos & Storage externo controlado & Checkpoints, manifests, datasets y, en la arquitectura objetivo, volúmenes y assets durables. \\
Entorno experimental & Google Colab, Drive y notebooks & Entrenamiento, evaluación, generación de checkpoints y evidencia reproducible. \\
\hline
\end{longtable}
}

% tab:arch-entidades-principales
{\footnotesize
\setlength{\LTleft}{0pt}
\setlength{\LTright}{0pt}
\setlength{\tabcolsep}{2pt}
\renewcommand{\arraystretch}{1.18}
\phantomsection\label{tab:arch-entidades-principales}
\begin{longtable}{>{\raggedright\arraybackslash}p{0.20\textwidth} >{\raggedright\arraybackslash}p{0.72\textwidth}}
\hline
\textbf{Entidad} & \textbf{Descripción y atributos clave} \\
\hline
\endfirsthead
\hline
\textbf{Entidad} & \textbf{Descripción y atributos clave} \\
\hline
\endhead
\hline
\endfoot
usuario & Cuenta del sistema. Atributos: id (PK), email, password\_hash, estado\_aprobacion, fecha\_registro. \\
rol / usuario\_rol & Roles ADMIN, DOCTOR y REVIEWER, asociados a usuarios mediante una tabla intermedia (relación muchos a muchos). \\
sujeto & Referencia de-identificada del paciente. Atributos: subject\_ref (PK). No almacena datos identificatorios. \\
estudio & Contenedor lógico de un caso. Atributos: study\_id (PK), subject\_ref (FK), creado\_por (FK usuario), fecha, metadata. \\
corrida (run) & Ejecución de análisis por plano. Atributos: run\_id (PK), study\_id (FK), model\_key (FK), plano, multiplanar\_run\_id, artifact\_hash, requested/effective\_inference\_mode, trace\_id, estado. \\
modelo / artifact & Modelo registrado y su artifact verificable. Atributos: model\_key (PK), plano, version, artifact\_hash, manifest\_uri, arquitectura. \\
medicion & Valor geométrico derivado de una máscara. Atributos: id (PK), run\_id (FK), estructura/label, plano, slice\_index, valor, unidad, confidence, estado\_revision. \\
revision & Decisión profesional sobre una corrida. Atributos: id (PK), run\_id (FK), usuario\_id (FK), estado (pendiente/aceptado/observado/rechazado/corregido), comentario, fecha. \\
asset & Archivo derivado servido bajo autorización. Atributos: id (PK), run\_id (FK), rol (input/mask/overlay/…), content\_type, hash, size, url. \\
slice\_catalogo (P10.5) & Catálogo de cortes del volumen para el visor. Atributos: id (PK), run\_id (FK), index, preview\_asset, has\_results. \\
evento\_auditoria & Registro de trazabilidad. Atributos: id (PK), usuario\_id (FK), entidad, accion, trace\_id, timestamp. \\
\hline
\end{longtable}
}

% tab:arch-red-seguridad
{\footnotesize
\setlength{\LTleft}{0pt}
\setlength{\LTright}{0pt}
\setlength{\tabcolsep}{2pt}
\renewcommand{\arraystretch}{1.18}
\phantomsection\label{tab:arch-red-seguridad}
\begin{longtable}{>{\raggedright\arraybackslash}p{0.16\textwidth} >{\raggedright\arraybackslash}p{0.16\textwidth} >{\raggedright\arraybackslash}p{0.24\textwidth} >{\raggedright\arraybackslash}p{0.36\textwidth}}
\hline
\textbf{Origen} & \textbf{Destino} & \textbf{Mecanismo} & \textbf{Propósito} \\
\hline
\endfirsthead
\hline
\textbf{Origen} & \textbf{Destino} & \textbf{Mecanismo} & \textbf{Propósito} \\
\hline
\endhead
\hline
\endfoot
Navegador & Frontend & HTTPS & Acceso a la interfaz. \\
Frontend & Backend & REST, JSON, JWT y X-Trace-Id & Operaciones de producto y trazabilidad. \\
Backend & AI Module & REST, JSON y X-Trace-Id & Inferencia, readiness, modelos y orquestación. \\
Backend & PostgreSQL & SQL/TLS según entorno & Persistencia transaccional. \\
AI Module & Storage & HTTPS o SDK autenticado & Descarga y verificación de models y manifests. \\
\hline
\end{longtable}
}

% tab:arch-artifacts-modelos
{\footnotesize
\setlength{\LTleft}{0pt}
\setlength{\LTright}{0pt}
\setlength{\tabcolsep}{2pt}
\renewcommand{\arraystretch}{1.18}
\phantomsection\label{tab:arch-artifacts-modelos}
\begin{longtable}{>{\raggedright\arraybackslash}p{0.18\textwidth} >{\raggedright\arraybackslash}p{0.34\textwidth} >{\raggedright\arraybackslash}p{0.40\textwidth}}
\hline
\textbf{Plano} & \textbf{Artifact principal} & \textbf{Función} \\
\hline
\endfirsthead
\hline
\textbf{Plano} & \textbf{Artifact principal} & \textbf{Función} \\
\hline
\endhead
\hline
\endfoot
Sagital & sagittal\_spider\_multiclass\_final\_best.pt & Segmentación multiclase de estructuras sagitales. \\
Axial & axial\_t2\_alkafri\_final\_best.pt & Segmentación axial según la codificación del dataset complementario. \\
\hline
\end{longtable}
}

% tab:arch-tecnologias
{\footnotesize
\setlength{\LTleft}{0pt}
\setlength{\LTright}{0pt}
\setlength{\tabcolsep}{2pt}
\renewcommand{\arraystretch}{1.18}
\phantomsection\label{tab:arch-tecnologias}
\begin{longtable}{>{\raggedright\arraybackslash}p{0.20\textwidth} >{\raggedright\arraybackslash}p{0.24\textwidth} >{\raggedright\arraybackslash}p{0.48\textwidth}}
\hline
\textbf{Tecnología} & \textbf{Rol} & \textbf{Justificación} \\
\hline
\endfirsthead
\hline
\textbf{Tecnología} & \textbf{Rol} & \textbf{Justificación} \\
\hline
\endhead
\hline
\endfoot
Spring Boot / Java & Backend de producto & Se utiliza como frontera de producto para concentrar autenticación, autorización, reglas de negocio, persistencia, auditoría y contratos públicos en un servicio independiente del runtime científico. Esta separación evita trasladar al Backend las dependencias propias del procesamiento y de los modelos de inteligencia artificial. \\
Python + FastAPI & Servicio IA & Se utiliza en el AI Module por su integración directa con PyTorch, SimpleITK y el código de preprocesamiento e inferencia. FastAPI permite exponer estas capacidades mediante una API HTTP interna consumida por el Backend, manteniendo desacoplado el runtime científico del resto del producto. \\
PyTorch & Entrenamiento e inferencia & Se mantiene como runtime porque los modelos y checkpoints del proyecto fueron entrenados, evaluados y registrados dentro de este ecosistema. Sustituirlo requeriría validar previamente la equivalencia de las salidas del nuevo runtime. \\
React + Vite + TypeScript & Frontend & Se utiliza para construir una interfaz web interactiva y modular. TypeScript aporta contratos tipados y React facilita la composición del visor, la worklist y los flujos de revisión. El Frontend consume exclusivamente el contrato público expuesto por el Backend. \\
PostgreSQL & Persistencia & Se utiliza por la necesidad de mantener integridad transaccional y relaciones consistentes entre usuarios, estudios, corridas, mediciones, revisiones y eventos de auditoría. \\
Docker & Portabilidad & Permite encapsular las dependencias y configuraciones de cada servicio, reproducir los entornos de ejecución y desplegar Backend y AI Module de manera desacoplada. \\
Google Colab + Drive & Experimentación & Se utilizan para entrenamiento, evaluación y conservación de la evidencia experimental y académica. No forman parte del runtime operativo del producto. \\
Google Cloud & Arquitectura objetivo & Se definió como plataforma objetivo por disponer de servicios administrados para cómputo, persistencia, almacenamiento y gestión operativa, manteniendo la separación entre componentes y permitiendo controlar los recursos asignados. Su implementación completa permanece planificada. \\
\hline
\end{longtable}
}

% tab:arch-alternativas
{\scriptsize
\setlength{\LTleft}{0pt}
\setlength{\LTright}{0pt}
\setlength{\tabcolsep}{2pt}
\renewcommand{\arraystretch}{1.18}
\phantomsection\label{tab:arch-alternativas}
\begin{longtable}{>{\raggedright\arraybackslash}p{0.18\textwidth} >{\raggedright\arraybackslash}p{0.23\textwidth} >{\raggedright\arraybackslash}p{0.31\textwidth} >{\raggedright\arraybackslash}p{0.20\textwidth}}
\hline
\textbf{Alternativa} & \textbf{Ventaja principal} & \textbf{Limitación o costo para el PFI} & \textbf{Decisión} \\
\hline
\endfirsthead
\hline
\textbf{Alternativa} & \textbf{Ventaja principal} & \textbf{Limitación o costo para el PFI} & \textbf{Decisión} \\
\hline
\endhead
\hline
\endfoot
Python + FastAPI para Backend y AI Module & Reduce la heterogeneidad tecnológica y permite utilizar un único lenguaje en los servicios del lado servidor. & No elimina la necesidad de separar las responsabilidades de producto del runtime científico y requeriría trasladar o reimplementar capacidades de autenticación, autorización, persistencia, auditoría y reglas de negocio ya establecidas en el Backend. & No adoptada. FastAPI se mantiene como servicio especializado de IA. \\
Java para Backend e inferencia & Unifica el stack del lado servidor y reduce la cantidad de lenguajes utilizados. & El pipeline de entrenamiento e inferencia, los checkpoints y las dependencias de procesamiento científico del proyecto se encuentran integrados con el ecosistema Python, PyTorch y SimpleITK. Migrarlos introduciría complejidad sin aportar una ventaja funcional al alcance actual. & No adoptada. Java se mantiene en el Backend de producto. \\
Spring Boot para Backend y Python + FastAPI para AI Module & Permite utilizar en cada componente el ecosistema más adecuado a su responsabilidad y mantener una frontera explícita entre producto e inteligencia artificial. & Obliga a mantener dos stacks tecnológicos y contratos de integración entre ambos servicios. & Adoptada. \\
ONNX Runtime para inferencia & Podría desacoplar la ejecución de los modelos del runtime original de PyTorch y facilitar alternativas de despliegue. & Requiere convertir los modelos y verificar que las salidas sean equivalentes a las obtenidas con los checkpoints actualmente evaluados. Esa equivalencia no fue validada en esta etapa. & Diferida. \\
\hline
\end{longtable}
}

% tab:arch-cloud-dimensionamiento
{\scriptsize
\setlength{\LTleft}{0pt}
\setlength{\LTright}{0pt}
\setlength{\tabcolsep}{2pt}
\renewcommand{\arraystretch}{1.18}
\phantomsection\label{tab:arch-cloud-dimensionamiento}
\begin{longtable}{>{\raggedright\arraybackslash}p{0.14\textwidth} >{\raggedright\arraybackslash}p{0.23\textwidth} >{\raggedright\arraybackslash}p{0.13\textwidth} >{\raggedright\arraybackslash}p{0.13\textwidth} >{\raggedright\arraybackslash}p{0.13\textwidth} >{\raggedright\arraybackslash}p{0.16\textwidth}}
\hline
\textbf{Componente} & \textbf{Escenario medido} & \textbf{CPU observada} & \textbf{RAM observada} & \textbf{GPU/VRAM} & \textbf{Criterio preliminar para despliegue} \\
\hline
\endfirsthead
\hline
\textbf{Componente} & \textbf{Escenario medido} & \textbf{CPU observada} & \textbf{RAM observada} & \textbf{GPU/VRAM} & \textbf{Criterio preliminar para despliegue} \\
\hline
\endhead
\hline
\endfoot
Frontend & Navegación y serving containerizado & p95: 1,98 \%; pico: 2,69 \% & p95: 14,68 MiB; pico: 14,80 MiB & No aplica & Componente de baja demanda computacional; apto para hosting estático o un contenedor de recursos mínimos. \\
Backend & Reposo y recorrido E2E & p95: 7,67 \%; pico: 33,40 \% & p95: 267,50 MiB; pico: 306,60 MiB & No aplica & La medición muestra una demanda moderada y ampliamente inferior a una vCPU durante el recorrido evaluado; se debe conservar margen de memoria para el runtime Java y variaciones de carga. \\
AI Module & Modelos cargados e inferencia real & p95: 346,62 \%; pico: 405,09 \% & p95: 801,20 MiB; pico: 940,80 MiB & GPU no utilizada & Es el principal consumidor de cómputo. El p95 equivale aproximadamente al uso concurrente de 3,47 vCPU y el pico a 4,05 vCPU en el entorno medido, por lo que el dimensionamiento inicial debe priorizar capacidad de CPU y margen de memoria. \\
PostgreSQL & Persistencia durante recorrido E2E & p95: 4,13 \%; pico: 30,81 \% & p95: 37,19 MiB; pico: 37,37 MiB & No aplica & La carga observada fue baja para un único recorrido E2E; corresponde utilizar una instancia administrada básica para el prototipo y reevaluarla bajo carga concurrente. \\
\hline
\end{longtable}
}

% tab:arch-interfaz-trazabilidad
{\scriptsize
\setlength{\LTleft}{0pt}
\setlength{\LTright}{0pt}
\setlength{\tabcolsep}{2pt}
\renewcommand{\arraystretch}{1.18}
\begin{longtable}{>{\raggedright\arraybackslash}p{0.22\textwidth} >{\raggedright\arraybackslash}p{0.30\textwidth} >{\raggedright\arraybackslash}p{0.16\textwidth} >{\raggedright\arraybackslash}p{0.24\textwidth}}
\caption{Relaci?n entre zonas de la interfaz, necesidades del relevamiento y requerimientos}\label{tab:arch-interfaz-trazabilidad}\\
\hline
\textbf{Zona de la interfaz} & \textbf{Función} & \textbf{Necesidad} & \textbf{Requerimiento} \\
\hline
\endfirsthead
\caption[]{Relaci?n entre zonas de la interfaz, necesidades del relevamiento y requerimientos (continuaci?n)}\\
\hline
\textbf{Zona de la interfaz} & \textbf{Función} & \textbf{Necesidad} & \textbf{Requerimiento} \\
\hline
\endhead
\hline
\endfoot
Worklist de estudios & Identificar casos disponibles, pendientes y revisados sin abrir cada uno. & N-07 & HU-02 \\
Alta de estudio y carga & Asociar el volumen a una referencia de sujeto seudonimizada e informar las series compatibles detectadas. & N-09, N-11 & HU-03, HU-04 \\
Visor con superposición conmutable & Comparar imagen original y máscara en el mismo encuadre. & N-01 & HU-06 \\
Selector de plano y navegación por cortes & Recorrer el volumen conservando la relación entre resultado y corte de origen. & N-03, N-12 & HU-07 \\
Panel de mediciones & Presentar cada valor junto con su estructura, plano y corte. & N-04 & HU-08 \\
Controles de revisión & Aceptar, observar, corregir o descartar cada resultado desde el punto en que se muestra. & N-02, N-05 & HU-09 \\
Preinforme editable & Reunir estructuras, mediciones, observaciones y estado de revisión en una salida modificable. & N-04, N-08 & HU-09 \\
Estados vacíos y degradados & Informar la ausencia de plano, secuencia o resultado sin completar la información faltante. & N-11 & HU-11 \\
Historial de corridas & Reabrir un estudio ya procesado sin repetir la inferencia. & N-08, N-10 & HU-10 \\
\hline
\end{longtable}
}